Determine todos \(r > 0\) tais que, sempre que \(f \colon \mathbb{R}^2 \to \mathbb{R}\) for uma função diferenciável tal que \(|\operatorname{grad}f(0,0)| = 1\) e \(|\operatorname{grad}f(u) - \operatorname{grad} f(v)| \le |u - v|\) para todos \(u,v \in \mathbb{R}^2\), então o máximo de \(f\) no disco \(\{u \in \mathbb{R}^2 : |u| \le r\}\) é atingido em exatamente um ponto.

(\(\operatorname{grad} f(u) = (\partial_1 f(u), \partial_2 f(u))\) é o vetor gradiente de \(f\) no ponto \(u\); para um vetor \(u = (a,b)\), \(|u| = \sqrt{a^2 + b^2}\).)
